\chapter{Capacitación contra Incendios y Emergencias}
\label{sec.cap7:capacitacion_emergencias}

La efectividad de las barreras tecnológicas y de ingeniería analizadas en los capítulos precedentes depende
fundamentalmente de la preparación y competencia operativa del factor humano. Frente a un conato de incendio, la
diferencia entre una contingencia sofocada en sus primeros treinta segundos y una catástrofe con colapso estructural
radica en el adiestramiento práctico del personal y en la fluidez de los protocolos de evacuación y autoprotección.

\section{Marco Regulatorio de la Capacitación Laboral en Argentina}

La formación en higiene y seguridad no constituye una recomendación optativa, sino un mandato legal taxativo:
\begin{itemize}
    \item Ley Nº 19.587 de Higiene y Seguridad en el Trabajo: En su Artículo 9, inciso k, establece la obligación del
          empleador de disponer de los medios necesarios para capacitar a todo el personal en materia de higiene y
          seguridad, prevención de siniestros y evacuación en caso de contingencias.
    \item Decreto Reglamentario 351/79 (Capítulo 21, Artículos 208 al 214): Dispone que los establecimientos deben
          ejecutar un programa anual de capacitación continua. Las actividades deben ser planificadas por profesionales
          habilitados, documentadas en actas formales firmadas por los trabajadores e integradas al legajo técnico de la
          empresa, a disposición permanente de las autoridades de control y de la Aseguradora de Riesgos del Trabajo (ART).
\end{itemize}

\section{Ítems Fundamentales a Considerar en una Capacitación contra Incendios}
\label{sec.cap7:items_capacitacion}

Al proyectar y dictar un programa formativo contra incendios en el ámbito industrial o electrónico, los responsables de
seguridad deben contemplar una serie de requerimientos técnicos, pedagógicos y operativos esenciales:

\subsection{1. Relevamiento y Diagnóstico Previo de los Riesgos Propios de la Instalación}
Una capacitación técnica no debe basarse en temarios genéricos y descontextualizados. El programa debe formularse a
partir de un relevamiento riguroso de la carga de fuego real del establecimiento (calculada conforme al Anexo VII del
Decreto 351/79), las materias primas manipuladas (disolventes, plásticos, aceites térmicos), la ubicación de tableros
eléctricos de potencia y las fuentes potenciales de ignición asociadas a las operaciones de la planta.

\subsection{2. Segmentación de Niveles Formativos según el Rol Ocupacional}
Los contenidos y la profundidad del adiestramiento deben estructurarse de manera diferenciada:
\begin{itemize}
    \item Nivel Básico (Todo el Personal): Enfocado en el reconocimiento temprano de anomalías, la activación del sistema
          de alarma, el uso correcto del extintor manual en fuegos incipientes y las pautas para una evacuación ordenada.
    \item Nivel Avanzado (Brigada de Emergencia): Destinado a los miembros del equipo de primera y segunda intervención,
          abarcando el tendido y operación de líneas de hidrantes, el uso de equipos de respiración autónoma (ERA), técnicas
          de rescate de víctimas atrapadas y coordinación con los cuarteles de bomberos públicos.
    \item Nivel Técnico Específico (Mantenimiento e Ingeniería): Focalizado en la prevención de fallas de arco eléctrico,
          bloqueo de fuentes de energía (procedimientos LOTO), protocolos de permisos de trabajo en caliente (soldadura y
          corte) y maniobras de corte de tableros en emergencias.
\end{itemize}

\subsection{3. Fundamentos Conceptuales del Fenómeno Ígneo}
Es indispensable que los trabajadores asimilen de forma intuitiva los conceptos fisicoquímicos básicos: el tetraedro del
fuego, las distintas clases de fuego normalizadas (Clases A, B, C, D y K) y el gravísimo peligro de aplicar un agente
inadecuado (e.g., el choque eléctrico letal al utilizar agua o espumas sobre tableros energizados Clase C, o la explosión
violenta de hidrógeno al arrojar agua sobre magnesio o sodio Clase D). Asimismo, se debe enfatizar la naturaleza
asfixiante y fulminante del monóxido de carbono y del cianuro de hidrógeno, desmitificando la idea de que el fuego avisa
mediante llamas visibles antes de resultar mortal.

\subsection{4. Entrenamiento Práctico y Destreza con Fuego Real}
La formación teórica de aula resulta inútil si los operarios no han manipulado previamente un extintor encendido. La
instrucción debe incluir prácticas de campo con fuego controlado sobre bateas de ensayo o simuladores electrónicos
certificados, permitiendo que el trabajador supere el bloqueo psicológico y la parálisis inducida por el estrés de las
llamas abiertas y el ruido de la combustión.

\subsection{5. Técnica Operativa Estandarizada de Extinción (Método PASAR / PASS)}
En la práctica se debe inculcar la secuencia técnica estandarizada para la descarga de extintores portátiles:
\begin{itemize}
    \item P - Pasador: Retirar el pasador de seguridad rompiendo el precinto plástico de control con un giro seco.
    \item A - Apuntar: Dirigir la tobera, manguera o difusor hacia la base de las llamas, donde se encuentra el
          combustible en descomposición, y nunca hacia la cúspide visible del fuego.
    \item S - Sujetar / Apretar: Presionar la palanca o gatillo de accionamiento con fuerza firme y continuada.
    \item A - Abanicar: Mover la tobera con un movimiento pendular de vaivén o zigzag de lado a lado sobre el frente de
          fuego, asegurando una cobertura homogénea del agente sobre la superficie ardiente.
    \item R - Retirarse: Una vez extinguido el foco, retroceder con marcha lenta mirando continuamente hacia el combustible,
          sin darle jamás la espalda al foco sofocado por el riesgo inminente de reignición imprevista.
\end{itemize}

\subsection{6. Posición Operativa y Factores Ambientales de Seguridad}
Durante el uso de extintores, el operario debe posicionarse a una distancia de seguridad inicial de $2\m$ a $3\m$ del
foco de fuego. Es terminantemente obligatorio ubicarse siempre a favor del viento (con el viento a la espalda),
impidiendo que las corrientes de convección térmica, el humo asfixiante y la nube de polvo químico sean impulsados contra
la cara y vías respiratorias del operador. Previo al acercamiento, se debe verificar que la aguja del manómetro se ubique
en la zona verde de presión operativa.

\subsection{7. Organización y Dinámica de la Brigada de Emergencias}
La capacitación debe formalizar la estructura operativa del plan de autoprotección:
\begin{itemize}
    \item Jefe de Emergencia / Comandante del Incidente: Máxima autoridad en el siniestro, responsable de la toma de
          decisiones estratégicas y del enlace directo con las fuerzas públicas exteriores (Bomberos y Policía).
    \item Brigadistas de Intervención: Personal entrenado que evalúa la magnitud del foco, corta la energía eléctrica del
          sector y procede a la extinción inicial mediante matafuegos o hidrantes.
    \item Guías de Evacuación: Encargados de verificar que todos los ocupantes de su piso o nave abandonen los puestos,
          asistir a personas con movilidad reducida y guiar el flujo hacia las salidas seguras.
    \item Equipo de Primeros Auxilios: Responsable de establecer el puesto médico avanzado y prestar atención inicial a
          heridos o intoxicados hasta el arribo de las ambulancias.
\end{itemize}

\subsection{8. Pautas de Supervivencia y Evacuación ante Atmósferas con Humo}
Los entrenamientos deben enfatizar el comportamiento humano frente al humo:
\begin{itemize}
    \item Desplazamiento a Ras del Suelo: Dado que el aire caliente y los gases tóxicos ascienden formando un manto
          superior que desciende progresivamente, los operarios deben desplazarse agachados o gateando. El estrato de aire
          comprendido en los primeros $30\unit{\centi\meter}$ a $50\unit{\centi\meter}$ sobre el piso mantiene los mayores
          niveles de oxígeno, menor temperatura y visibilidad suficiente para reconocer la ruta de escape.
    \item Verificación Térmica de Puertas: Antes de abrir una puerta cerrada en una vía de escape, se debe tocar la hoja y
          el picaporte metálico con el dorso de la mano. Si la superficie presenta calor intenso, la puerta no debe abrirse,
          ya que el ingreso repentino de aire exterior puede provocar una explosión de humo o llamarada súbita.
    \item Prohibición Absoluta de Ascensores: En caso de emergencia queda taxativamente vedado el empleo de elevadores o
          montacargas, debiendo utilizarse exclusivamente cajas de escaleras presurizadas o salidas de emergencia.
\end{itemize}

\subsection{9. Punto de Encuentro Seguro y Control de Asistencia (Roll Call)}
El plan de evacuación debe designar un Punto de Encuentro ubicado al aire libre, a una distancia
segura del edificio superior a la altura máxima de su fachada, fuera del alcance del calor radiante y de posibles caídas
de mampostería o cables de media tensión. Una vez concentrado el personal, los guías de evacuación deben ejecutar de forma
inmediata el recuento de personas confrontando con el registro diario de ingreso. Esta información es
crítica para notificar fehacientemente al jefe de bomberos si existen personas atrapadas en sectores específicos.

\subsection{10. Planificación y Evaluación de Simulacros Periódicos}
Los simulacros de evacuación deben planificarse con una periodicidad mínima de dos veces al año. Se recomienda iniciar con
ejercicios programados con aviso previo para asimilar las rutas, transitando luego hacia simulacros imprevistos que
simulen contingencias reales (e.g., una salida de emergencia clausurada por humo). Todo simulacro debe ser cronometrado,
evaluando los tiempos de desocupación total, y culminar en una reunión de análisis técnico con el
Comité de Seguridad para registrar desvíos y actualizar los planes de autoprotección.

\section{Equipamiento de Protección Personal (EPP) y Brigada}

La dotación de EPP constituye la última línea de defensa física. Para la intervención en emergencias, la selección técnica
debe responder a la matriz expuesta en la Tabla \ref{tab.cap7:epp}.

\begin{table}[!ht]
    \centering
    \resizebox{\textwidth}{!}{
    \begin{tabular}{|>{\raggedright\arraybackslash}p{4.2cm}|>{\raggedright\arraybackslash}p{4.8cm}|>{\raggedright\arraybackslash}p{7cm}|}
        \hline
        \textbf{Riesgo Específico} & \textbf{EPP Normalizado Requerido} & \textbf{Fundamento Físico y Fisiológico} \\ \hline
        Fogonazo y Arco Eléctrico en Tableros &
        Traje completo de fibras de aramida (Nomex/Kevlar), capucha y visor de policarbonato con protección UV/IR. &
        Las fibras poliaramidas no funden ni gotean ante temperaturas de plasma de hasta $6000\Celsius$, protegiendo la
        piel contra quemaduras de tercer grado por radiación extrema. \\ \hline
        Combate Directo de Incendios en Naves &
        Traje estructural de bombero multicapa (Norma NFPA 1971), guantes de cuero ignífugo y botas de caucho vulcanizado. &
        La barrera de humedad y el forro térmico retardan la transferencia de calor por conducción y convección, aislando
        al brigadista del contacto con vapores calientes y agua hirviente. \\ \hline
        Atmósferas Contaminadas con Humo Tóxico &
        Equipos de Respiración Autónoma (ERA) de presión positiva con cilindro de aire comprimido a $300\,\mathrm{bar}$. &
        Dado que el CO y el HCN no pueden filtrarse con mascarillas comunes de carbón activo en deficiencia de oxígeno, la
        presión positiva en la máscara garantiza que, ante un eventual desajuste perimetral, el aire fluya hacia el exterior
        impidiendo la entrada de tóxicos. \\ \hline
    \end{tabular}
    }
    \caption{Matriz técnica de Equipos de Protección Personal (EPP) para brigadas de incendio.}
    \label{tab.cap7:epp}
\end{table}
